\documentclass[12pt]{article}
\usepackage{amsmath,amssymb,amsfonts,amsthm}
\usepackage[margin=1in]{geometry}

\begin{document}

\begin{center}
{\Large\bfseries Chapter 9, Problem 18}\\[6pt]
Byron \& Fuller, \textit{Mathematics of Classical and Quantum Physics}
\end{center}

\bigskip

\textbf{Problem.} Show that if $k(x,y) = |\psi(x)|^2\psi^*(y)$ (a rank-1 kernel... actually $k(x,y) = \psi(x)\psi^*(y)$), then the Fredholm method reduces to $D(\lambda) = 1 - \lambda\int |\psi(t)|^2\,dt$, $N(x,y;\lambda) = k(x,y)$, and all higher terms vanish. Show this leads to a solution of $f = g + \lambda Kf$ in agreement with Eq.~(8.57).

\bigskip
\textbf{Solution.}

\subsection*{Fredholm series for a rank-1 kernel}

Let $k(x,y) = \psi(x)\psi^*(y)$. This is a rank-1 (separable) kernel.

\textbf{Fredholm determinant $D(\lambda)$:}

The first-order term is:
\[
d_1 = \int k(x,x)\,dx = \int |\psi(x)|^2\,dx = \|\psi\|^2.
\]

The second-order term involves:
\[
d_2 = \int\!\!\int \begin{vmatrix} k(x_1,x_1) & k(x_1,x_2) \\ k(x_2,x_1) & k(x_2,x_2) \end{vmatrix} dx_1\,dx_2.
\]

Computing the determinant:
\begin{align}
&k(x_1,x_1)k(x_2,x_2) - k(x_1,x_2)k(x_2,x_1) \\
&= |\psi(x_1)|^2|\psi(x_2)|^2 - \psi(x_1)\psi^*(x_2)\psi(x_2)\psi^*(x_1) \\
&= |\psi(x_1)|^2|\psi(x_2)|^2 - |\psi(x_1)|^2|\psi(x_2)|^2 = 0.
\end{align}

So $d_2 = 0$. Similarly, all $n \times n$ determinants with $n \ge 2$ vanish because the matrix $K_n$ with entries $k(x_i, x_j) = \psi(x_i)\psi^*(x_j)$ has rank 1, so its determinant is zero for $n \ge 2$.

Therefore:
\[
\boxed{D(\lambda) = 1 - \lambda\|\psi\|^2}.
\]

\textbf{Fredholm numerator $N(x,y;\lambda)$:}

$N_0(x,y) = k(x,y) = \psi(x)\psi^*(y)$.

For $n \ge 1$, $N_n$ involves an $(n+1) \times (n+1)$ determinant where the first row/column involves $(x,y)$ and the remaining $n$ rows/columns form a rank-1 matrix. The full $(n+1) \times (n+1)$ matrix also has rank at most 2 (from the outer product structure), and in fact all the $n \times n$ minors needed vanish for $n \ge 1$ by the same rank argument.

Therefore:
\[
\boxed{N(x,y;\lambda) = k(x,y) = \psi(x)\psi^*(y)}.
\]

\subsection*{Solution of the integral equation}

The resolvent is:
\[
R(x,y;\lambda) = \frac{N(x,y;\lambda)}{D(\lambda)} = \frac{\psi(x)\psi^*(y)}{1 - \lambda\|\psi\|^2}.
\]

The solution to $f = g + \lambda Kf$ is:
\[
f(x) = g(x) + \lambda\int R(x,y;\lambda)\,g(y)\,dy = g(x) + \frac{\lambda\,\psi(x)}{1 - \lambda\|\psi\|^2}\int \psi^*(y)\,g(y)\,dy.
\]

This can be written as:
\[
\boxed{f(x) = g(x) + \frac{\lambda\,(\psi, g)}{1 - \lambda\|\psi\|^2}\,\psi(x)},
\]
which agrees with the standard result for separable kernels (Eq.~8.57). The solution exists for all $\lambda \neq 1/\|\psi\|^2$. \hfill $\blacksquare$

\end{document}
